% !TEX program = xelatex
\documentclass[a4paper,10pt]{article}
\usepackage[UTF8,fontset=none]{ctex}
\usepackage{xeCJK}
\setCJKmainfont{Songti SC}
\setCJKsansfont{STHeiti}
\setCJKmonofont{Songti SC}
\usepackage[margin=2.2cm]{geometry}
\usepackage{tikz}
\usetikzlibrary{shapes.geometric, arrows.meta, positioning, calc, fit}
\usepackage{amsmath, amssymb}
\usepackage{longtable}
\usepackage{array}
\usepackage{booktabs}
\usepackage{multirow}
\usepackage{enumitem}
\usepackage{fancyhdr}
\usepackage{titlesec}
\usepackage{hyperref}
\usepackage{xcolor}
\usepackage{tabularx}

\hypersetup{
  colorlinks=true,
  linkcolor=blue!60!black,
  urlcolor=blue!60!black,
  bookmarksnumbered=true
}

% 页眉页脚
\pagestyle{fancy}
\fancyhf{}
\fancyhead[L]{\small evaluatePolynomial~---~函数设计文档}
\fancyhead[R]{\small 内部公开}
\fancyfoot[C]{\thepage}
\renewcommand{\headrulewidth}{0.4pt}

% 层级编号
\setcounter{secnumdepth}{3}
\setcounter{tocdepth}{3}

% 表格列宽辅助
\newcolumntype{L}[1]{>{\raggedright\arraybackslash}p{#1}}
\newcolumntype{C}[1]{>{\centering\arraybackslash}p{#1}}

% ============================================================
\begin{document}

% ==================== 封面表格 ====================
\thispagestyle{empty}

\begin{center}
{\CJKfamily{zhsong}\bfseries\Large 函数设计文档}
\end{center}

\vspace{8pt}

\begin{center}
\renewcommand{\arraystretch}{1.5}
\begin{tabular}{|L{3.2cm}|L{4.5cm}|L{3.2cm}|L{4.5cm}|}
\hline
\textbf{函数英文名} & evaluatePolynomial &
\textbf{函数中文名} & 多项式求值 \\
\hline
\textbf{函数类型} & 元件 (Element) &
\textbf{版本} & 2.0 \\
\hline
\textbf{作者} & 许庆 &
\textbf{日期} & 2026-06-26 \\
\hline
\textbf{联系方式} & \multicolumn{3}{l|}{18612426814~/~qingxu@tsinghua.edu.cn} \\
\hline
\textbf{密级} & 内部公开 &
\textbf{AI辅助} & Cursor Agent (Claude Opus 4.6) \\
\hline
\end{tabular}
\end{center}

\vspace{12pt}

\begin{center}
\renewcommand{\arraystretch}{1.4}
\begin{tabular}{|C{1.2cm}|C{2cm}|L{5cm}|C{2cm}|C{2cm}|}
\hline
\multicolumn{5}{|c|}{\textbf{更改记录}} \\
\hline
\textbf{版本} & \textbf{日期} & \textbf{更改说明} & \textbf{作者} & \textbf{辅助AI} \\
\hline
2.0 & 2026-06-26 & 初始版本 & 许庆 & Claude Opus 4.6 \\
\hline
\end{tabular}
\end{center}

\newpage

% ==================== 目录 ====================
\tableofcontents
\newpage

% ============================================================
%  1. 函数需求
% ============================================================
\section{函数需求}

\subsection{功能概述}

\texttt{evaluatePolynomial} 使用 Horner 法（秦九韶算法）计算三次多项式在给定点 $x$ 处的值：
\[
  y = c_0 + c_1 x + c_2 x^2 + c_3 x^3
\]
Horner 法将其改写为嵌套乘法形式以减少乘法运算次数并提高数值稳定性。

\subsection{输入/输出/返回值}

\renewcommand{\arraystretch}{1.3}
\begin{longtable}{|C{1.5cm}|L{3.5cm}|C{3cm}|L{7cm}|}
\hline
\textbf{类别} & \textbf{名称} & \textbf{类型} & \textbf{说明} \\
\hline
\endfirsthead
\hline
\textbf{类别} & \textbf{名称} & \textbf{类型} & \textbf{说明} \\
\hline
\endhead
输入 & \texttt{coeffs} & const PolynomialCoeffs\& & 多项式系数 $c[0..3]$ \\
\hline
输入 & \texttt{x} & double & 自变量值 \\
\hline
返回值 & --- & double & 多项式在 $x$ 处的函数值 \\
\hline
\end{longtable}

% ============================================================
%  1.2 算法设计
% ============================================================
\subsection{算法设计}

三次多项式的标准形式为：
\begin{equation}
  y = c_0 + c_1 x + c_2 x^2 + c_3 x^3
\end{equation}

直接展开需要 6 次乘法和 3 次加法。Horner 法将其改写为嵌套形式：
\begin{equation}
  y = c_0 + x \cdot \big(c_1 + x \cdot (c_2 + x \cdot c_3)\big)
\end{equation}

仅需 3 次乘法和 3 次加法，计算效率提升 50\%。同时由于从高阶向低阶逐步累加，
中间结果的量级递减，有利于减少浮点舍入误差。

\textbf{计算步骤}：
\begin{enumerate}
  \item $t_1 = c_2 + x \cdot c_3$
  \item $t_2 = c_1 + x \cdot t_1$
  \item $y = c_0 + x \cdot t_2$
\end{enumerate}

实际实现中合并为单行表达式：
\[
  \texttt{result} = \texttt{coeffs.c[0]} + x \times (\texttt{coeffs.c[1]} + x \times (\texttt{coeffs.c[2]} + x \times \texttt{coeffs.c[3]}))
\]

% ============================================================
%  1.3 程序流程图
% ============================================================
\subsection{程序流程图}

\begin{center}
\begin{tikzpicture}[
  node distance=1.0cm,
  startstop/.style={rectangle, rounded corners, minimum width=3.5cm, minimum height=0.7cm,
    text centered, draw=black, fill=blue!8, font=\small},
  process/.style={rectangle, minimum width=5.5cm, minimum height=0.7cm,
    text centered, draw=black, fill=orange!8, font=\small},
  arrow/.style={-{Stealth[length=2.5mm]}, thick}
]

\node (start) [startstop] {开始};
\node (input) [process, below=of start] {读入 coeffs, x};
\node (calc) [process, below=of input] {result = c[0] + x*(c[1] + x*(c[2] + x*c[3]))};
\node (output) [process, below=of calc] {返回 result};
\node (stop) [startstop, below=of output] {结束};

\draw [arrow] (start) -- (input);
\draw [arrow] (input) -- (calc);
\draw [arrow] (calc) -- (output);
\draw [arrow] (output) -- (stop);

\end{tikzpicture}
\end{center}

% ============================================================
%  1.4 函数接口
% ============================================================
\subsection{函数接口}

\begin{verbatim}
double evaluatePolynomial(const PolynomialCoeffs& coeffs, const double x);
\end{verbatim}

% ============================================================
%  1.5 输入输出模块图
% ============================================================
\subsection{输入输出模块图}

\begin{center}
\begin{tikzpicture}[
  box/.style={rectangle, draw=black, thick, minimum width=3.5cm, minimum height=0.6cm,
    text centered, fill=yellow!10, font=\small},
  core/.style={rectangle, draw=black, very thick, minimum width=5cm, minimum height=1.5cm,
    text centered, fill=orange!10, font=\normalsize\bfseries, rounded corners=3pt},
  arrow/.style={-{Stealth[length=2.5mm]}, thick}
]

\node (core) [core] {evaluatePolynomial};

\node (in1) [box, left=3cm of core, yshift=0.4cm] {coeffs (PolynomialCoeffs)};
\node (in2) [box, left=3cm of core, yshift=-0.4cm] {x (double)};

\node (out1) [box, right=3cm of core] {y (double)};

\draw [arrow] (in1) -- (core);
\draw [arrow] (in2) -- (core);
\draw [arrow] (core) -- (out1);

\node [above=0.1cm of in1, font=\small\bfseries] {输入 (Input)};
\node [above=0.1cm of out1, font=\small\bfseries] {输出 (Output)};

\end{tikzpicture}
\end{center}

% ============================================================
%  1.6 函数诸元表
% ============================================================
\subsection{函数诸元表}

\renewcommand{\arraystretch}{1.3}
\begin{longtable}{|L{3.8cm}|L{11.5cm}|}
\hline
\textbf{属性} & \textbf{说明} \\
\hline
\endfirsthead
\hline
\textbf{属性} & \textbf{说明} \\
\hline
\endhead
函数英文名 & \texttt{evaluatePolynomial} \\
\hline
函数中文名 & 多项式求值 \\
\hline
函数类型 & 元件 (Element) \\
\hline
版本 & 2.0 \\
\hline
源文件 & \texttt{src/cpp/evaluatePolynomial/evaluatePolynomial.cpp} \\
\hline
头文件 & \texttt{include/cpp/evaluatePolynomial/evaluatePolynomial.hpp} \\
\hline
命名空间 & \texttt{waypoint\_follow} \\
\hline
输入数量 & 2（coeffs, x） \\
\hline
输出数量 & 1（返回值 double） \\
\hline
参数数量 & 0 \\
\hline
状态数量 & 0（无状态，纯函数） \\
\hline
下级函数数量 & 0 \\
\hline
动态内存分配 & 无 \\
\hline
外部依赖 & 无（仅算术运算） \\
\hline
算法类别 & 多项式计算 / Horner 法 \\
\hline
\end{longtable}

% ============================================================
%  1.7 函数架构
% ============================================================
\subsection{函数架构}

无下级函数调用。\texttt{evaluatePolynomial} 为叶节点元件，仅包含算术运算。

% ============================================================
%  1.8 数据结构定义
% ============================================================
\subsection{数据结构定义}

引用 \texttt{waypointFollowTypes.hpp} 中定义的结构体：

\subsubsection{PolynomialCoeffs~---~三次多项式系数}
\begin{longtable}{|L{3.5cm}|C{2cm}|L{9.5cm}|}
\hline
\textbf{成员} & \textbf{类型} & \textbf{说明} \\
\hline
\texttt{c[4]} & double[4] & $y = c[0] + c[1]\,x + c[2]\,x^2 + c[3]\,x^3$ \\
\hline
\end{longtable}

% ============================================================
%  1.9 测试用例
% ============================================================
\subsection{测试用例}

\subsubsection{正常测试用例}

{\small
\begin{longtable}{|C{0.6cm}|L{2.2cm}|L{4.5cm}|L{3.5cm}|L{3.8cm}|}
\hline
\textbf{编号} & \textbf{场景} & \textbf{输入} & \textbf{预期输出} & \textbf{验证准则} \\
\hline
\endfirsthead
\hline
\textbf{编号} & \textbf{场景} & \textbf{输入} & \textbf{预期输出} & \textbf{验证准则} \\
\hline
\endhead
N-1 &
零点求值 &
$\mathbf{c}{=}[1,2,3,4]$, $x{=}0$ &
$y = 1.0$ &
$y = c_0$（高阶项消失） \\
\hline
N-2 &
$x{=}1$ 全系数 &
$\mathbf{c}{=}[1,2,3,4]$, $x{=}1$ &
$y = 10.0$ &
$1{+}2{+}3{+}4 = 10$ \\
\hline
N-3 &
负 $x$ 值 &
$\mathbf{c}{=}[1,2,3,4]$, $x{=}{-}1$ &
$y = -2.0$ &
$1{-}2{+}3{-}4 = -2$ \\
\hline
N-4 &
仅常数项 &
$\mathbf{c}{=}[5,0,0,0]$, $x{=}100$ &
$y = 5.0$ &
零系数不影响结果 \\
\hline
N-5 &
小数值 $x$ &
$\mathbf{c}{=}[0,1,0,0]$, $x{=}0.5$ &
$y = 0.5$ &
$y = c_1 \cdot x = 0.5$ \\
\hline
\end{longtable}
}

\vspace{6pt}
{\small
\begin{longtable}{|L{2.5cm}|L{13cm}|}
\hline
\multicolumn{2}{|c|}{\textbf{正常测试用例符号说明}} \\
\hline
\textbf{符号} & \textbf{含义} \\
\hline
$\mathbf{c}$ & 多项式系数数组 coeffs.c[0..3] \\
\hline
$x$ & 自变量值 \\
\hline
$y$ & 返回值（多项式函数值） \\
\hline
$c_i$ & 第 $i$ 阶系数，即 coeffs.c[$i$] \\
\hline
\end{longtable}
}

\subsubsection{异常测试用例}

{\small
\begin{longtable}{|C{0.6cm}|L{2.0cm}|L{4.2cm}|L{4.0cm}|L{3.8cm}|}
\hline
\textbf{编号} & \textbf{场景} & \textbf{输入} & \textbf{预期输出/行为} & \textbf{验证准则} \\
\hline
\endfirsthead
\hline
\textbf{编号} & \textbf{场景} & \textbf{输入} & \textbf{预期输出/行为} & \textbf{验证准则} \\
\hline
\endhead
A-1 &
$x$ 为 NaN &
$\mathbf{c}{=}[1,2,3,4]$, $x{=}\text{NaN}$ &
$y = \text{NaN}$ &
NaN 参与运算传播 \\
\hline
A-2 &
$x$ 为 $+\infty$ &
$\mathbf{c}{=}[1,2,3,4]$, $x{=}+\infty$ &
$y = +\infty$ &
$c_3 > 0$ 时结果为 $+\infty$ \\
\hline
A-3 &
$x$ 为 $-\infty$ &
$\mathbf{c}{=}[1,2,3,4]$, $x{=}-\infty$ &
$y = -\infty$ &
$c_3 > 0$, 奇数次幂主导 \\
\hline
A-4 &
系数含 NaN &
$\mathbf{c}{=}[1,\text{NaN},3,4]$, $x{=}2$ &
$y = \text{NaN}$ &
NaN 系数传播 \\
\hline
A-5 &
系数含 $\infty$ &
$\mathbf{c}{=}[+\infty,0,0,0]$, $x{=}1$ &
$y = +\infty$ &
$\infty$ 系数传播 \\
\hline
A-6 &
$\infty - \infty$ 抵消 &
$\mathbf{c}{=}[+\infty,0,0,-\infty]$, $x{=}1$ &
$y = \text{NaN}$ &
$\infty + (-\infty) = $ NaN \\
\hline
A-7 &
极大 $x$ 值 &
$\mathbf{c}{=}[0,0,0,1]$, $x{=}10^{100}$ &
$y = 10^{300}$ &
未溢出（$< $ DBL\_MAX） \\
\hline
A-8 &
溢出场景 &
$\mathbf{c}{=}[0,0,0,1]$, $x{=}10^{103}$ &
$y = +\infty$（溢出） &
$10^{309} > $ DBL\_MAX \\
\hline
A-9 &
全零系数 &
$\mathbf{c}{=}[0,0,0,0]$, $x{=}10^{308}$ &
$y = 0.0$ &
精确零结果 \\
\hline
A-10 &
subnormal $x$ &
$\mathbf{c}{=}[0,1,0,0]$, $x{=}5\!\times\!10^{-324}$ &
$y \approx 5\!\times\!10^{-324}$ &
极小值正常处理 \\
\hline
\end{longtable}
}

\vspace{6pt}
{\small
\begin{longtable}{|L{2.5cm}|L{13cm}|}
\hline
\multicolumn{2}{|c|}{\textbf{异常测试用例符号说明}} \\
\hline
\textbf{符号} & \textbf{含义} \\
\hline
$\mathbf{c}$ & 多项式系数数组 coeffs.c[0..3] \\
\hline
$x$ & 自变量值 \\
\hline
$y$ & 返回值（多项式函数值） \\
\hline
NaN & IEEE 754 非数值，任何含 NaN 的算术运算结果为 NaN \\
\hline
$+\infty$/$-\infty$ & IEEE 754 正/负无穷大 \\
\hline
DBL\_MAX & double 最大有限值 $\approx 1.7\!\times\!10^{308}$ \\
\hline
subnormal & 非规格化浮点数（极接近零） \\
\hline
\end{longtable}
}

\end{document}
